\scanpage{55}
\mspage{54}

majorées par $F$ s'appellent les \uline{strates} de $\deleted{\unsure{E}}F$. (Donc, \deleted{\unsure{de|des}} strate de $F$ est une multistrate)

L'ens des strates de $F$ \deleted{s'appelle} sera noté $\overline F$ : \[ \begin{aligned} \deleted{\unsure{\overline F\subset\mathcal M|F\in\mathcal M}} &\qquad \unsure{\mathcal M|\mathscr M}\subset\unsure{\mathcal F|\mathscr F}\\ F\longrightarrow\overline F &\qquad \mathcal F\longrightarrow\unsure{\mathcal P|\mathscr P}(\mathcal M). \end{aligned} \]

\uline{$\unsure{\mathcal C|\mathscr C|\mathcal E}$ 2} Toute figure est la borne supérieure de ses strates : \[ F=\underset{\deleted{\unsure{A}}X\in\overline F}{\mathrm{Sup}}\ X \]

\uline{Corollaire} 1 L'application $\deleted{X}F\longrightarrow\overline{\deleted{\unsure{C|G}}F}$ \[ \mathcal F\longrightarrow\mathcal P(\mathcal M) \] est \unsure{injection|injective}, et induit un iso d'\unsure{ens} ordonnés de $\mathcal F$ avec l'ens $\unsure{\Phi|Q|\varphi}\subset\mathcal P(\mathcal M)$ formé des $\overline F$.

Cet iso. transforme $\underset{i}{\mathrm{Sup}}$ en $\bigcup_i$

\uline{Corollaire 2} L'ens $\Phi$ des parties de $\mathcal M$ \unsure{satisfait|satisfaisant|satisfont} les deux conditions

\deleted{a) $\Phi\ne\varnothing$}

% slanted
\margin{\deleted{\unsure{Ceci|Ceci ?}}} \deleted{\unsure{b)}}a) Si $A\in\Phi$, et $(A_i)_{i\in\unsure{I|\Phi}}$ est une famille \unsure{di'|d'}éléments de $\Phi$, avec $A_i\subset A$, alors $\bigcup_iA_i\in\Phi$.

b) $\forall X\in\mathcal M$, il existe un plus petit $A=A_X$ tel que $X\in\unsure{A|A_x|A_X}$ (a fortiori, $\Phi$ recouvre $\mathcal M$).
